\documentclass[aspectratio=169]{beamer}
\usetheme{Madrid}
\usecolortheme{whale}
\usepackage[utf8]{inputenc}
\usepackage[spanish]{babel}
\usepackage{amsmath,amssymb}
\usepackage{graphicx}
\usepackage{booktabs}
\usepackage{listings}
\usepackage{xcolor}

\definecolor{codegreen}{rgb}{0,0.6,0}
\definecolor{codegray}{rgb}{0.5,0.5,0.5}
\definecolor{codepurple}{rgb}{0.58,0,0.82}
\definecolor{backcolour}{rgb}{0.95,0.95,0.92}

\lstdefinestyle{pythonstyle}{
    backgroundcolor=\color{backcolour},
    commentstyle=\color{codegreen},
    keywordstyle=\color{blue},
    numberstyle=\tiny\color{codegray},
    stringstyle=\color{codepurple},
    basicstyle=\ttfamily\footnotesize,
    breakatwhitespace=false,
    breaklines=true,
    captionpos=b,
    keepspaces=true,
    showspaces=false,
    showstringspaces=false,
    showtabs=false,
    tabsize=2
}

\lstset{style=pythonstyle}

\title{Regresión Lineal Simple}
\subtitle{Conceptos y Estimación}
\author{Métodos de Análisis de Datos 1}
\institute{Universidad Nacional del Sur}
\date{Unidad 3 - Clase 14}

\begin{document}

\frame{\titlepage}

\begin{frame}{Contenidos}
\tableofcontents
\end{frame}

%==============================================================================
\section{Introducción al Modelado}
%==============================================================================

\begin{frame}{¿Qué es un modelo estadístico?}
\begin{block}{Definición}
Un modelo estadístico es una representación matemática de un fenómeno del mundo real que usa datos observados.
\end{block}

\vspace{0.5cm}
\textbf{Objetivos principales:}
\begin{itemize}
    \item \textbf{Descripción:} Entender cómo se relacionan las variables
    \item \textbf{Predicción:} Estimar valores futuros de una variable
    \item \textbf{Inferencia:} Determinar qué variables son importantes
    \item \textbf{Control:} Ajustar por factores confusores
\end{itemize}
\end{frame}

\begin{frame}{Modelos de regresión}
La \textbf{regresión} modela la relación entre:
\begin{itemize}
    \item Variable \textbf{respuesta} $Y$ (dependiente, outcome)
    \item Variable(s) \textbf{predictora(s)} $X$ (independiente(s), explicativa(s))
\end{itemize}

\vspace{0.5cm}
\begin{table}
\centering
\begin{tabular}{ll}
\toprule
\textbf{Tipo} & \textbf{Características} \\
\midrule
Lineal simple & 1 predictor, relación lineal \\
Lineal múltiple & $\geq 2$ predictores, relación lineal \\
Polinomial & 1 predictor, relación no lineal \\
Logística & Variable respuesta binaria \\
\bottomrule
\end{tabular}
\end{table}

\vspace{0.3cm}
\textbf{Esta clase:} Regresión lineal simple
\end{frame}

%==============================================================================
\section{Modelo de Regresión Lineal Simple}
%==============================================================================

\begin{frame}{Definición del modelo}
\begin{block}{Modelo de regresión lineal simple}
\[
Y_i = \beta_0 + \beta_1 X_i + \epsilon_i, \quad i = 1, 2, \ldots, n
\]
\end{block}

donde:
\begin{itemize}
    \item $Y_i$: valor observado de la respuesta
    \item $X_i$: valor del predictor
    \item $\beta_0$: \textbf{intercepto} (ordenada al origen)
    \item $\beta_1$: \textbf{pendiente}
    \item $\epsilon_i$: \textbf{error aleatorio}, $\epsilon_i \sim N(0, \sigma^2)$ independientes
    \item $n$: tamaño de la muestra
\end{itemize}
\end{frame}

\begin{frame}{Interpretación gráfica}
\begin{center}
\includegraphics[width=0.7\textwidth]{regresion_simple.png}
\end{center}

\begin{itemize}
    \item La línea representa $E(Y|X) = \beta_0 + \beta_1 X$
    \item Los puntos son los datos observados $(X_i, Y_i)$
    \item Las distancias verticales son los residuos $\epsilon_i$
\end{itemize}
\end{frame}

\begin{frame}{Interpretación de los parámetros}
\begin{block}{Intercepto $\beta_0$}
Valor esperado de $Y$ cuando $X = 0$
\end{block}

\begin{block}{Pendiente $\beta_1$}
Cambio esperado en $Y$ cuando $X$ aumenta en 1 unidad
\end{block}

\vspace{0.5cm}
\textbf{Ejemplo:} Precio de casas vs metros cuadrados
\[
Y = 50 + 1.2 X
\]
\begin{itemize}
    \item $\beta_0 = 50$: precio base (50 mil \$)
    \item $\beta_1 = 1.2$: cada m$^2$ adicional suma 1.2 mil \$ (1200 \$)
\end{itemize}
\end{frame}

%==============================================================================
\section{Supuestos del Modelo}
%==============================================================================

\begin{frame}{Supuestos fundamentales}
Para que la inferencia sea válida:

\begin{enumerate}
    \item \textbf{Linealidad:} La relación entre $X$ e $Y$ es lineal
    \item \textbf{Independencia:} Las observaciones son independientes
    \item \textbf{Homocedasticidad:} Varianza constante de los errores
    \[
    \text{Var}(\epsilon_i) = \sigma^2 \text{ para todo } i
    \]
    \item \textbf{Normalidad:} Los errores siguen una distribución normal
    \[
    \epsilon_i \sim N(0, \sigma^2)
    \]
\end{enumerate}

\vspace{0.3cm}
\textbf{Nota:} Los veremos en detalle en la clase de diagnóstico
\end{frame}

\begin{frame}{¿Qué pasa si se violan los supuestos?}
\begin{itemize}
    \item \textbf{Linealidad:} Predicciones sesgadas, bajo $R^2$
    \item \textbf{Independencia:} Errores estándar incorrectos (muy optimistas)
    \item \textbf{Homocedasticidad:} Errores estándar incorrectos, intervalos inválidos
    \item \textbf{Normalidad:} Tests e intervalos menos confiables (especialmente con $n$ pequeño)
\end{itemize}

\vspace{0.5cm}
\begin{alertblock}{Importante}
La próxima clase veremos cómo \textbf{diagnosticar} y \textbf{corregir} estas violaciones
\end{alertblock}
\end{frame}

%==============================================================================
\section{Estimación por Mínimos Cuadrados}
%==============================================================================

\begin{frame}{Método de Mínimos Cuadrados Ordinarios (MCO)}
\textbf{Objetivo:} Encontrar la recta que mejor se ajusta a los datos

\vspace{0.3cm}
\textbf{Criterio:} Minimizar la suma de los cuadrados de los residuos
\[
\text{SCR} = \sum_{i=1}^{n} (Y_i - \hat{Y}_i)^2 = \sum_{i=1}^{n} (Y_i - \hat{\beta}_0 - \hat{\beta}_1 X_i)^2
\]

donde $\hat{Y}_i = \hat{\beta}_0 + \hat{\beta}_1 X_i$ es el valor predicho.

\vspace{0.5cm}
\begin{center}
\includegraphics[width=0.55\textwidth]{residuos.png}
\end{center}
\end{frame}

\begin{frame}{Fórmulas de MCO}
Las estimaciones que minimizan SCR son:

\begin{block}{Pendiente}
\[
\hat{\beta}_1 = \frac{\sum_{i=1}^{n} (X_i - \bar{X})(Y_i - \bar{Y})}{\sum_{i=1}^{n} (X_i - \bar{X})^2} = \frac{S_{XY}}{S_{XX}}
\]
\end{block}

\begin{block}{Intercepto}
\[
\hat{\beta}_0 = \bar{Y} - \hat{\beta}_1 \bar{X}
\]
\end{block}

donde:
\begin{itemize}
    \item $\bar{X}, \bar{Y}$: medias muestrales
    \item $S_{XY} = \sum (X_i - \bar{X})(Y_i - \bar{Y})$: suma de productos cruzados
    \item $S_{XX} = \sum (X_i - \bar{X})^2$: suma de cuadrados de $X$
\end{itemize}
\end{frame}

\begin{frame}{Ejemplo numérico}
Datos: Altura ($X$, cm) vs Peso ($Y$, kg) de 5 personas

\begin{table}
\centering
\begin{tabular}{ccccc}
\toprule
$X_i$ & $Y_i$ & $X_i - \bar{X}$ & $Y_i - \bar{Y}$ & $(X_i-\bar{X})(Y_i-\bar{Y})$ \\
\midrule
160 & 55 & -10 & -10 & 100 \\
165 & 60 & -5 & -5 & 25 \\
170 & 65 & 0 & 0 & 0 \\
175 & 70 & 5 & 5 & 25 \\
180 & 75 & 10 & 10 & 100 \\
\midrule
$\bar{X}=170$ & $\bar{Y}=65$ & & & $S_{XY}=250$ \\
\bottomrule
\end{tabular}
\end{table}

\vspace{0.3cm}
$S_{XX} = 10^2 + 5^2 + 0^2 + 5^2 + 10^2 = 250$

\vspace{0.3cm}
$\hat{\beta}_1 = \frac{250}{250} = 1.0$, \quad $\hat{\beta}_0 = 65 - 1.0 \times 170 = -105$

\vspace{0.3cm}
Modelo: $\hat{Y} = -105 + 1.0 X$
\end{frame}

\begin{frame}{Propiedades de MCO}
Los estimadores por MCO tienen propiedades deseables:

\begin{itemize}
    \item \textbf{Insesgados:} $E(\hat{\beta}_0) = \beta_0$, $E(\hat{\beta}_1) = \beta_1$
    \item \textbf{Consistentes:} Convergen al valor verdadero cuando $n \to \infty$
    \item \textbf{BLUE (Best Linear Unbiased Estimator):} 
    \begin{itemize}
        \item Tienen la menor varianza entre todos los estimadores lineales insesgados
        \item Teorema de Gauss-Markov
    \end{itemize}
\end{itemize}

\vspace{0.5cm}
\begin{alertblock}{Importante}
Estas propiedades se cumplen \textbf{bajo los supuestos del modelo}
\end{alertblock}
\end{frame}

%==============================================================================
\section{Coeficiente de Determinación $R^2$}
%==============================================================================

\begin{frame}{Descomposición de la variabilidad}
La variabilidad total de $Y$ se descompone en:

\[
\underbrace{\sum_{i=1}^{n} (Y_i - \bar{Y})^2}_{\text{SCT (Total)}} = \underbrace{\sum_{i=1}^{n} (\hat{Y}_i - \bar{Y})^2}_{\text{SCE (Explicada)}} + \underbrace{\sum_{i=1}^{n} (Y_i - \hat{Y}_i)^2}_{\text{SCR (Residual)}}
\]

\vspace{0.5cm}
\begin{itemize}
    \item \textbf{SCT:} Variabilidad total de $Y$ (sin modelo)
    \item \textbf{SCE:} Variabilidad explicada por el modelo
    \item \textbf{SCR:} Variabilidad no explicada (error)
\end{itemize}
\end{frame}

\begin{frame}{Coeficiente de determinación $R^2$}
\begin{block}{Definición}
\[
R^2 = \frac{\text{SCE}}{\text{SCT}} = 1 - \frac{\text{SCR}}{\text{SCT}}
\]
\end{block}

\textbf{Interpretación:} Proporción de variabilidad de $Y$ explicada por $X$

\vspace{0.5cm}
\textbf{Propiedades:}
\begin{itemize}
    \item $0 \leq R^2 \leq 1$
    \item $R^2 = 0$: el modelo no explica nada
    \item $R^2 = 1$: ajuste perfecto
    \item En regresión simple: $R^2 = r_{XY}^2$ (cuadrado de la correlación)
\end{itemize}

\vspace{0.3cm}
\begin{alertblock}{Advertencia}
$R^2$ alto no implica causalidad ni garantiza que el modelo sea adecuado
\end{alertblock}
\end{frame}

\begin{frame}{Ejemplo de $R^2$}
Continuando con el ejemplo altura-peso:

\vspace{0.3cm}
Modelo: $\hat{Y} = -105 + 1.0 X$

\begin{table}
\centering
\begin{tabular}{cccc}
\toprule
$Y_i$ & $\hat{Y}_i$ & $(Y_i - \bar{Y})^2$ & $(Y_i - \hat{Y}_i)^2$ \\
\midrule
55 & 55 & 100 & 0 \\
60 & 60 & 25 & 0 \\
65 & 65 & 0 & 0 \\
70 & 70 & 25 & 0 \\
75 & 75 & 100 & 0 \\
\midrule
& & SCT = 250 & SCR = 0 \\
\bottomrule
\end{tabular}
\end{table}

\vspace{0.3cm}
$R^2 = 1 - \frac{0}{250} = 1.0$ (ajuste perfecto!)
\end{frame}

%==============================================================================
\section{Implementación en Python}
%==============================================================================

\begin{frame}[fragile]{Regresión lineal con statsmodels}
\begin{lstlisting}[language=Python]
import pandas as pd
import statsmodels.api as sm
from statsmodels.formula.api import ols

# Datos
data = pd.DataFrame({
    'altura': [160, 165, 170, 175, 180],
    'peso': [55, 60, 65, 70, 75]
})

# Ajustar modelo
modelo = ols('peso ~ altura', data=data).fit()

# Resumen completo
print(modelo.summary())

# Coeficientes
print(f"Intercepto: {modelo.params['Intercept']:.2f}")
print(f"Pendiente: {modelo.params['altura']:.4f}")
print(f"R^2: {modelo.rsquared:.4f}")
\end{lstlisting}
\end{frame}

\begin{frame}[fragile]{Visualización}
\begin{lstlisting}[language=Python]
import matplotlib.pyplot as plt

# Grafico de dispersion y recta ajustada
plt.figure(figsize=(8, 5))
plt.scatter(data['altura'], data['peso'], 
            label='Datos', s=80, alpha=0.6)
plt.plot(data['altura'], modelo.fittedvalues, 
         'r-', linewidth=2, label='Ajuste')
plt.xlabel('Altura (cm)')
plt.ylabel('Peso (kg)')
plt.title('Regresion Lineal Simple')
plt.legend()
plt.grid(True, alpha=0.3)
plt.show()
\end{lstlisting}
\end{frame}

\begin{frame}[fragile]{Predicción}
\begin{lstlisting}[language=Python]
# Predecir para nuevos valores
nuevos_datos = pd.DataFrame({'altura': [172, 178, 168]})

# Prediccion puntual
predicciones = modelo.predict(nuevos_datos)
print("Predicciones:", predicciones.values)

# Prediccion con intervalo de confianza
pred_intervalo = modelo.get_prediction(nuevos_datos)
intervalos = pred_intervalo.summary_frame(alpha=0.05)
print(intervalos)

# Columnas: mean, mean_se, mean_ci_lower, mean_ci_upper,
#           obs_ci_lower, obs_ci_upper
\end{lstlisting}
\end{frame}

%==============================================================================
\section{Implementación en R}
%==============================================================================

\begin{frame}[fragile]{Regresión lineal en R}
\begin{lstlisting}[language=R]
# Datos
altura <- c(160, 165, 170, 175, 180)
peso <- c(55, 60, 65, 70, 75)
data <- data.frame(altura, peso)

# Ajustar modelo
modelo <- lm(peso ~ altura, data = data)

# Resumen
summary(modelo)

# Coeficientes
coef(modelo)

# R^2
summary(modelo)$r.squared
\end{lstlisting}
\end{frame}

\begin{frame}[fragile]{Visualización en R}
\begin{lstlisting}[language=R]
# Grafico basico
plot(data$altura, data$peso, 
     xlab = "Altura (cm)", ylab = "Peso (kg)",
     main = "Regresion Lineal Simple",
     pch = 19, col = "blue")
abline(modelo, col = "red", lwd = 2)

# Con ggplot2
library(ggplot2)
ggplot(data, aes(x = altura, y = peso)) +
  geom_point(size = 3, alpha = 0.6) +
  geom_smooth(method = "lm", se = TRUE, 
              color = "red", fill = "pink") +
  theme_minimal() +
  labs(title = "Regresion Lineal Simple")
\end{lstlisting}
\end{frame}

\begin{frame}[fragile]{Predicción en R}
\begin{lstlisting}[language=R]
# Nuevos datos
nuevos <- data.frame(altura = c(172, 178, 168))

# Prediccion puntual
predict(modelo, nuevos)

# Intervalo de confianza para la media
predict(modelo, nuevos, interval = "confidence")

# Intervalo de prediccion para nuevas observaciones
predict(modelo, nuevos, interval = "prediction")

# El intervalo de prediccion es siempre mas ancho
# que el de confianza
\end{lstlisting}
\end{frame}

%==============================================================================
\section{Resumen}
%==============================================================================

\begin{frame}{Resumen de la clase}
\textbf{Aprendimos:}
\begin{itemize}
    \item El modelo de regresión lineal simple: $Y = \beta_0 + \beta_1 X + \epsilon$
    \item Interpretación de intercepto y pendiente
    \item Los 4 supuestos fundamentales (L.I.H.N.)
    \item Estimación por mínimos cuadrados (MCO)
    \item Coeficiente de determinación $R^2$
    \item Implementación en Python (statsmodels) y R (lm)
\end{itemize}

\vspace{0.5cm}
\textbf{Próxima clase:}
\begin{itemize}
    \item Inferencia sobre los coeficientes (tests e intervalos)
    \item Predicción (intervalos de confianza vs predicción)
    \item ANOVA del modelo de regresión
\end{itemize}
\end{frame}

\end{document}
